\subsection{Streaming Data}
The main data source are single column comma separated values files which are read by a stream of c++ library "fstreams". The stream is an independent object which is created in the main file with a corresponding file path. \\ \\
\lstinputlisting[
  language=C,
  caption={Object which handles the streaming of data.},
  label={src:linregressionalgorithm},
  firstnumber=1,
  firstline=1,
  lastline=36
]{content/codebase/streamingData.hpp}
Rather than loading complete dataset batches into memory, incoming sensor data is  processed line per line.
Data ingestion is encapsulated within a custom stream handler class (\texttt{streamData}), which wraps standard C++ file input streams (\texttt{std::ifstream}). The boolean flag (\texttt{bool next}) and (\texttt{bool hasNext}) are used to check if there are still some values left in the file which can be read by the handler.
To eliminate potential data crashes, data races and runtime memory reallocation overhead, all operational state buffers are pre-allocated on the stack during configuration.
Finally, the resulting principal component scores are accumulated in a dynamic buffer and written to disk (\texttt{output.txt}) upon successful completion of the algorithm execution.